En este informe se estudian dos problemas: el
ordenamiento de arreglos unidimensionales y la multiplicación de matrices
cuadradas. Para el problema de ordenamiento se analizan experimentalmente cuatro algoritmos de ordenamiento ($std::sort$, $PatienceSort$, $QuickSort$ y $MergeSort$) y para el de matrices se analizan dos algoritmos ($Naive$ y $Strassen$). El objetivo principal es analizar el comportamiento práctico tanto para tiempo de ejecución (ms) como para consumo de memoria (KB), pudiendo identificar y discutir las diferencias que pueden presentarse con lo dictado en la teoría.
Para lograr este objetivo se utilizan casos de prueba generados aleatoriamente, con diferentes tamaños y elementos, para luego visualizar mediante gráficos los resultados cuando estas variables cambian.